% =====================================================================
% Introduction (with integrated literature review) -- DRAFT
% Usage: % =====================================================================
% Introduction (with integrated literature review) -- DRAFT
% Usage: % =====================================================================
% Introduction (with integrated literature review) -- DRAFT
% Usage: % =====================================================================
% Introduction (with integrated literature review) -- DRAFT
% Usage: \input{introduction_draft} in place of the current
%        "Introduction", "Main Contributions" and
%        "Solution Approaches in the Literature" blocks.
% Requires: natbib (\citet/\citep), xcolor, tabularx, booktabs.
% Red text = numbers or claims to confirm before submission.
% =====================================================================

\section{Introduction}
\label{sec:introduction}

The inventory routing problem (IRP) arises in vendor-managed inventory
systems, in which a supplier monitors the stock levels of its customers and
takes responsibility for replenishing them. Over a multi-period planning
horizon, the supplier jointly decides when each customer is visited, how much
product is delivered at each visit, and how the vehicles are routed, subject
to customer storage capacities, vehicle capacities, and the requirement that
no customer runs out of stock. Integrating these decisions, rather than
fixing replenishment plans first and routing them afterwards, can produce
substantial savings \citep{archetti2016inventory}. However, the integration
couples a vehicle routing problem in every period with inventory dynamics
that link the periods together. As a result, the IRP remains one of the most
challenging problems in distribution logistics
\citep{moin2007overview,bertazzi2012inventory,coelho2014thirty,ekcsi2026inventory}.

Most of the IRP literature minimises the total cost, defined as the sum of
routing and inventory holding costs. In many applications, however, holding
costs are difficult to estimate or are negligible compared with transport
costs, and practitioners evaluate distribution plans by the transport cost
per unit of product delivered. This key performance indicator, known as the
\emph{logistic ratio}, has long been used in the distribution of industrial
gases \citep{bell1983improving,benoist2011randomized,singh2015incremental}.
It was also the objective of the ROADEF/EURO Challenge 2016, proposed by
Air Liquide \citep{andre2020roadef}. \citet{archetti2017logistic} formalised
the IRP with logistic ratio (IRP-LR) and showed that solutions that minimise
total cost can be poor in terms of the logistic ratio. One reason is the
finite-horizon effect: under a cost objective there is little incentive to
deliver product near the end of the horizon if it will only be consumed
afterwards. The logistic ratio, in contrast, rewards every unit delivered
and therefore produces replenishment plans that differ markedly from
cost-minimising ones.

The IRP-LR is harder to solve than the classical IRP. Its objective is the
ratio of two decision-dependent quantities: the routing cost in the numerator
and the total delivered quantity in the denominator. A move that shortens the
routes can therefore worsen the objective if it also reduces the quantity
delivered, and a longer route can be preferable if it allows a larger
delivery. Exact approaches build on parametric fractional programming
\citep{dinkelbach1967}, which replaces the ratio with a sequence of linear
objectives. \citet{archetti2017logistic} and \citet{archetti2019exact}
developed exact algorithms for the IRP-LR along these lines, and
\citet{baldacci2020fractional} proposed a general exact framework for vehicle
routing problems with fractional objectives. These methods provide optimal
solutions and bounds, but their computational requirements restrict them to
instances with at most \textcolor{red}{50} customers. For larger instances, to the best of our
knowledge, the only heuristics evaluated on the IRP-LR benchmark are the
iterated local search (ILS) and simulated annealing (SA) algorithms of
\citet{alvarez2018iterated}. The algorithms developed for the ROADEF/EURO
Challenge 2016 address a richer industrial variant with a ratio objective,
using heuristic branch-cut-and-price \citep{absi2020branch}, local search
combined with mathematical programming \citep{su2020matheuristic},
fixed-sequence reoptimisation of quantities and timing
\citep{he2020fixedsequence}, and hyper-heuristic sequence selection
\citep{kheiri2020heuristic}. More recently, ratio objectives have been
studied in a multi-source setting solved by ILS \citep{alonsopecina2024iterated}
and in a multi-product setting with lateral transshipment and uncertain
demand \citep{man2025logistics}. Compared with the classical IRP, for which
a wide range of exact and heuristic methods exists, the heuristic literature
on the IRP-LR remains sparse.

\subsection{Related literature and research gaps}
\label{sec:intro_literature}

For the classical IRP, exact methods based on branch-and-cut
\citep{archetti2007branch,coelho2013branch}, branch-price-and-cut
\citep{desaulniers2016branch}, and their matheuristic extensions
\citep{skaalnes2023branch} solve instances of moderate size. Large instances
are handled by heuristics, which are reviewed by \citet{sofianopoulou2021review}
and \citet{ekcsi2026inventory}. We organise the discussion around three
design choices that are particularly important under a logistic-ratio
objective. Table~\ref{tab:literature_irplr} summarises the studies most
closely related to this work.

\paragraph{Where delivery quantities are decided.}
Efficient local search is available for routing decisions, whereas delivery
quantities are naturally determined by mathematical programming. Most IRP
heuristics therefore decompose the problem: a heuristic explores visit
schedules and routes, and a mathematical program sets or repairs the
quantities \citep{iswari2026extending}. Examples include the tabu search
combined with mixed-integer programming of \citet{archetti2012hybrid}, the
matheuristic of \citet{archetti2017matheuristic}, the variable MIP
neighbourhood descent of \citet{coelho2020variable}, the ILS matheuristic of
\citet{lagana2024iterated}, and route-improving matheuristics for related
production routing problems \citep{vadseth2023multi}. In the hybrid genetic
search (HGS) of \citet{castro2026hybrid}, the routing component builds on the
HGS for the capacitated vehicle routing problem (CVRP)
\citep{vidal2022hybrid}, while the quantities are obtained by solving a
minimum-cost flow problem with the network simplex algorithm.
\citet{zhao2025large} extend HGS with a large neighbourhood that
simultaneously changes visit periods, insertion positions, and delivery
quantities across the horizon. Decomposition works well when holding costs
make quantity decisions secondary to routing. Under the logistic ratio,
however, the change in delivered quantity caused by a routing move is part
of the value of that move. A move that is evaluated with quantities held
fixed can appear unattractive even though it becomes improving once the
quantities are adjusted, and solving a mathematical program for every
candidate move is too expensive within a local search. This suggests
evaluating each move by its combined effect on routing cost and delivered
quantity. It also suggests using an exact optimisation step for quantities
only occasionally, as \citet{he2020fixedsequence} showed with fixed-sequence
reoptimisation.

\paragraph{How infeasibility is handled.}
Allowing a search to visit infeasible solutions and penalising the violation
in the objective is a well-established strategy in vehicle routing. It is a
central component of HGS \citep{vidal2022hybrid}, and in inventory routing
it has been used to relax vehicle capacity
\citep{archetti2012hybrid,castro2026hybrid} and customer stockouts
\citep{zhao2025large}. In the IRP, the stockout constraints make the
feasible region particularly fragmented. Removing a single visit usually
causes a stockout in some later period, so a search restricted to feasible
solutions cannot easily move between replenishment plans with different
visit patterns. Penalising stockouts addresses this, but under a ratio
objective the penalty weight is difficult to set. The logistic ratio is
measured in cost per unit and varies widely across instances, so a fixed
penalty is either too weak, and the search remains infeasible, or too
strong, and the search behaves as if infeasibility were forbidden. Lagrangian
relaxation \citep{fisher1981lagrangian} offers a principled alternative: the
stockout constraint is dualised, and the multiplier is updated with a
subgradient-type rule that scales the penalty to the current gap from the
incumbent. To the best of our knowledge, this has not previously been used
to guide a local search for the IRP-LR.

\paragraph{Saturated vehicles.}
Because the logistic ratio rewards large deliveries, vehicles in good
IRP-LR solutions are typically fully loaded. Any move that inserts a
customer into such a vehicle is then infeasible if the existing quantities
are kept, so large parts of the inter-route and inter-period neighbourhoods
are effectively unavailable. In the CVRP, HGS addresses the corresponding
issue by allowing capacity violations at a penalty \citep{vidal2022hybrid}.
In the IRP-LR, there is a feasible alternative. The deliveries already on a
full vehicle often include quantities that are not needed to prevent future
stockouts, and this slack can be reassigned to the inserted customer without
violating any constraint. Exploiting this slack keeps the search within
capacity while giving access to the moves that matter most for the
logistic ratio.

Together, these observations identify the gap addressed in this paper:
there is no scalable heuristic for the IRP-LR that (i) evaluates routing and
quantity decisions jointly within each move, (ii) explores
inventory-infeasible solutions in a controlled manner that adapts to the
ratio objective, and (iii) does not depend on a mathematical programming
solver.

\begin{table*}[!htbp]
\centering
\caption{Studies most relevant to the IRP-LR and to the proposed solution
methodology.}
\label{tab:literature_irplr}
\scriptsize
\renewcommand{\arraystretch}{1.18}
\setlength{\tabcolsep}{3pt}
\begin{tabularx}{\textwidth}{
>{\raggedright\arraybackslash}p{2.3cm}
>{\raggedright\arraybackslash}p{2.6cm}
>{\raggedright\arraybackslash}p{2.8cm}
>{\raggedright\arraybackslash}X
>{\raggedright\arraybackslash}X}
\toprule
\textbf{Study} & \textbf{Problem / objective} & \textbf{Solution approach}
& \textbf{Key search characteristic} & \textbf{Relevance to this work} \\
\midrule
\citet{archetti2017logistic} & IRP; routing cost / quantity delivered
  & Exact fractional optimisation & Joint routing and replenishment
  & Introduces the IRP-LR \\
\citet{alvarez2018iterated} & IRP including logistic ratio
  & ILS and SA & Neighbourhood search over routes and replenishment
  & Benchmark heuristic for the IRP-LR \\
\citet{archetti2019exact} & IRP-LR & Iterative exact method
  & Sequence of IRPs with linear objectives
  & Optimal solutions for small instances \\
\citet{baldacci2020fractional} & VRP with fractional objective
  & Exact method & Primal and dual bounding procedures
  & General fractional routing methodology \\
\citet{absi2020branch} & Rich IRP; unit distribution cost
  & Heuristic branch-cut-and-price
  & Integrated routing, quantity, and scheduling
  & Large-scale industrial ratio IRP \\
\citet{su2020matheuristic} & ROADEF/EURO IRP & Local search with MIP/LP
  & Mathematical programming sets timing and quantities
  & Hybrid local-search design \\
\citet{he2020fixedsequence} & ROADEF/EURO IRP
  & Fixed-sequence matheuristic
  & Reoptimises quantities and timing for fixed visit sequences
  & Separation of route structure and quantity optimisation \\
\citet{alonsopecina2024iterated} & Operating cost / quantity delivered
  & ILS & Neighbourhood search for a multi-source IRP
  & Recent ILS with a ratio objective \\
\citet{man2025logistics} & Multi-product logistics-ratio IRP
  & Robust fractional optimisation & Uncertainty and lateral transshipment
  & Extension of logistics-ratio optimisation \\
\citet{castro2026hybrid} & Classical IRP & HGS
  & Route search with min-cost-flow quantity optimisation
  & State-of-the-art decomposition-based search \\
\citet{zhao2025large} & Classical and stockout-penalised IRP
  & HGS with large neighbourhood
  & Cross-period route/quantity moves; penalised stockouts
  & Joint moves and controlled infeasibility \\
\midrule
\textbf{This paper} & IRP-LR & Multi-start ILS (solver-free)
  & Joint route--quantity moves; vehicle slack; Lagrangian stockout penalty
  & Scalable heuristic for the IRP-LR \\
\bottomrule
\end{tabularx}
\end{table*}

\subsection{Contributions}
\label{sec:intro_contributions}

This paper proposes a multi-start iterated local search (MSILS) for the
IRP-LR. The algorithm does not require a mathematical programming solver.
Its main contributions are as follows.

\begin{enumerate}
  \item \textbf{Joint evaluation of route and quantity decisions.} Four of
  the five local search operators update delivery quantities as part of
  each move, using closed-form rules that respect depot supply, vehicle
  capacity, and customer storage capacity. Every candidate move is
  therefore evaluated by its actual effect on the logistic ratio, that is,
  on routing cost and delivered quantity together.

  \item \textbf{Vehicle slack.} We introduce a mechanism that reassigns
  quantity not needed to prevent future stockouts from the existing
  deliveries of a fully loaded vehicle to a newly inserted customer. The
  mechanism is motivated by the capacity-violation strategy of HGS
  \citep{vidal2022hybrid}, but it never violates vehicle capacity. It makes
  available the inter-route and inter-period moves that would otherwise be
  rejected, and it creates no new stockouts.

  \item \textbf{Lagrangian-guided exploration of infeasible solutions.}
  The zero-stockout constraint is dualised, and the local search solves a
  sequence of relaxed problems in which the Lagrange multiplier is updated
  by a subgradient-type rule based on the gap from the incumbent. The
  penalty therefore scales automatically to the magnitude of the logistic
  ratio of each instance, allowing the search to cross infeasible regions
  and then return to feasibility.

  \item \textbf{Problem-specific construction, rebalancing, and
  perturbation.} Initial solutions are built by a randomised constructive
  heuristic that prioritises customers by stockout urgency and distance
  from the depot. A rebalance procedure redistributes quantities over a
  fixed schedule whenever the incumbent improves, and an adaptive
  perturbation increases its strength as the number of non-improving
  iterations grows.

  \item \textbf{Computational evidence.} Experiments on
  \textcolor{red}{1128} benchmark instances, comprising \textcolor{red}{978}
  instances with 5 to 50 customers and \textcolor{red}{150} instances with
  50 to 200 customers, show that the MSILS improves on the ILS and SA of
  \citet{alvarez2018iterated} by on average \textcolor{red}{6.30\%} and
  \textcolor{red}{8.00\%}, respectively, on the large instances, with larger
  gains as the instance size grows. On small instances, it finds solutions
  within \textcolor{red}{X\%} of the optimal values reported by
  \citet{archetti2019exact}. A component analysis quantifies the
  contribution of each mechanism, and we also examine whether embedding the
  HGS of \citet{vidal2022hybrid} as a routing post-optimiser brings further
  improvement. The instances are publicly available
  \textcolor{red}{(and the code will be released upon acceptance)}.
\end{enumerate}

The remainder of the paper is organised as follows.
Section~\ref{sec:problem_description} defines the IRP-LR and the
Lagrangian-relaxed problem used by the algorithm. Section~\ref{sec: alg}
describes the MSILS. Section~\ref{sec: computational experiements} reports
the computational experiments, and Section~\textcolor{red}{\ref{sec:conclusion}}
concludes and outlines directions for future research.
 in place of the current
%        "Introduction", "Main Contributions" and
%        "Solution Approaches in the Literature" blocks.
% Requires: natbib (\citet/\citep), xcolor, tabularx, booktabs.
% Red text = numbers or claims to confirm before submission.
% =====================================================================

\section{Introduction}
\label{sec:introduction}

The inventory routing problem (IRP) arises in vendor-managed inventory
systems, in which a supplier monitors the stock levels of its customers and
takes responsibility for replenishing them. Over a multi-period planning
horizon, the supplier jointly decides when each customer is visited, how much
product is delivered at each visit, and how the vehicles are routed, subject
to customer storage capacities, vehicle capacities, and the requirement that
no customer runs out of stock. Integrating these decisions, rather than
fixing replenishment plans first and routing them afterwards, can produce
substantial savings \citep{archetti2016inventory}. However, the integration
couples a vehicle routing problem in every period with inventory dynamics
that link the periods together. As a result, the IRP remains one of the most
challenging problems in distribution logistics
\citep{moin2007overview,bertazzi2012inventory,coelho2014thirty,ekcsi2026inventory}.

Most of the IRP literature minimises the total cost, defined as the sum of
routing and inventory holding costs. In many applications, however, holding
costs are difficult to estimate or are negligible compared with transport
costs, and practitioners evaluate distribution plans by the transport cost
per unit of product delivered. This key performance indicator, known as the
\emph{logistic ratio}, has long been used in the distribution of industrial
gases \citep{bell1983improving,benoist2011randomized,singh2015incremental}.
It was also the objective of the ROADEF/EURO Challenge 2016, proposed by
Air Liquide \citep{andre2020roadef}. \citet{archetti2017logistic} formalised
the IRP with logistic ratio (IRP-LR) and showed that solutions that minimise
total cost can be poor in terms of the logistic ratio. One reason is the
finite-horizon effect: under a cost objective there is little incentive to
deliver product near the end of the horizon if it will only be consumed
afterwards. The logistic ratio, in contrast, rewards every unit delivered
and therefore produces replenishment plans that differ markedly from
cost-minimising ones.

The IRP-LR is harder to solve than the classical IRP. Its objective is the
ratio of two decision-dependent quantities: the routing cost in the numerator
and the total delivered quantity in the denominator. A move that shortens the
routes can therefore worsen the objective if it also reduces the quantity
delivered, and a longer route can be preferable if it allows a larger
delivery. Exact approaches build on parametric fractional programming
\citep{dinkelbach1967}, which replaces the ratio with a sequence of linear
objectives. \citet{archetti2017logistic} and \citet{archetti2019exact}
developed exact algorithms for the IRP-LR along these lines, and
\citet{baldacci2020fractional} proposed a general exact framework for vehicle
routing problems with fractional objectives. These methods provide optimal
solutions and bounds, but their computational requirements restrict them to
instances with at most \textcolor{red}{50} customers. For larger instances, to the best of our
knowledge, the only heuristics evaluated on the IRP-LR benchmark are the
iterated local search (ILS) and simulated annealing (SA) algorithms of
\citet{alvarez2018iterated}. The algorithms developed for the ROADEF/EURO
Challenge 2016 address a richer industrial variant with a ratio objective,
using heuristic branch-cut-and-price \citep{absi2020branch}, local search
combined with mathematical programming \citep{su2020matheuristic},
fixed-sequence reoptimisation of quantities and timing
\citep{he2020fixedsequence}, and hyper-heuristic sequence selection
\citep{kheiri2020heuristic}. More recently, ratio objectives have been
studied in a multi-source setting solved by ILS \citep{alonsopecina2024iterated}
and in a multi-product setting with lateral transshipment and uncertain
demand \citep{man2025logistics}. Compared with the classical IRP, for which
a wide range of exact and heuristic methods exists, the heuristic literature
on the IRP-LR remains sparse.

\subsection{Related literature and research gaps}
\label{sec:intro_literature}

For the classical IRP, exact methods based on branch-and-cut
\citep{archetti2007branch,coelho2013branch}, branch-price-and-cut
\citep{desaulniers2016branch}, and their matheuristic extensions
\citep{skaalnes2023branch} solve instances of moderate size. Large instances
are handled by heuristics, which are reviewed by \citet{sofianopoulou2021review}
and \citet{ekcsi2026inventory}. We organise the discussion around three
design choices that are particularly important under a logistic-ratio
objective. Table~\ref{tab:literature_irplr} summarises the studies most
closely related to this work.

\paragraph{Where delivery quantities are decided.}
Efficient local search is available for routing decisions, whereas delivery
quantities are naturally determined by mathematical programming. Most IRP
heuristics therefore decompose the problem: a heuristic explores visit
schedules and routes, and a mathematical program sets or repairs the
quantities \citep{iswari2026extending}. Examples include the tabu search
combined with mixed-integer programming of \citet{archetti2012hybrid}, the
matheuristic of \citet{archetti2017matheuristic}, the variable MIP
neighbourhood descent of \citet{coelho2020variable}, the ILS matheuristic of
\citet{lagana2024iterated}, and route-improving matheuristics for related
production routing problems \citep{vadseth2023multi}. In the hybrid genetic
search (HGS) of \citet{castro2026hybrid}, the routing component builds on the
HGS for the capacitated vehicle routing problem (CVRP)
\citep{vidal2022hybrid}, while the quantities are obtained by solving a
minimum-cost flow problem with the network simplex algorithm.
\citet{zhao2025large} extend HGS with a large neighbourhood that
simultaneously changes visit periods, insertion positions, and delivery
quantities across the horizon. Decomposition works well when holding costs
make quantity decisions secondary to routing. Under the logistic ratio,
however, the change in delivered quantity caused by a routing move is part
of the value of that move. A move that is evaluated with quantities held
fixed can appear unattractive even though it becomes improving once the
quantities are adjusted, and solving a mathematical program for every
candidate move is too expensive within a local search. This suggests
evaluating each move by its combined effect on routing cost and delivered
quantity. It also suggests using an exact optimisation step for quantities
only occasionally, as \citet{he2020fixedsequence} showed with fixed-sequence
reoptimisation.

\paragraph{How infeasibility is handled.}
Allowing a search to visit infeasible solutions and penalising the violation
in the objective is a well-established strategy in vehicle routing. It is a
central component of HGS \citep{vidal2022hybrid}, and in inventory routing
it has been used to relax vehicle capacity
\citep{archetti2012hybrid,castro2026hybrid} and customer stockouts
\citep{zhao2025large}. In the IRP, the stockout constraints make the
feasible region particularly fragmented. Removing a single visit usually
causes a stockout in some later period, so a search restricted to feasible
solutions cannot easily move between replenishment plans with different
visit patterns. Penalising stockouts addresses this, but under a ratio
objective the penalty weight is difficult to set. The logistic ratio is
measured in cost per unit and varies widely across instances, so a fixed
penalty is either too weak, and the search remains infeasible, or too
strong, and the search behaves as if infeasibility were forbidden. Lagrangian
relaxation \citep{fisher1981lagrangian} offers a principled alternative: the
stockout constraint is dualised, and the multiplier is updated with a
subgradient-type rule that scales the penalty to the current gap from the
incumbent. To the best of our knowledge, this has not previously been used
to guide a local search for the IRP-LR.

\paragraph{Saturated vehicles.}
Because the logistic ratio rewards large deliveries, vehicles in good
IRP-LR solutions are typically fully loaded. Any move that inserts a
customer into such a vehicle is then infeasible if the existing quantities
are kept, so large parts of the inter-route and inter-period neighbourhoods
are effectively unavailable. In the CVRP, HGS addresses the corresponding
issue by allowing capacity violations at a penalty \citep{vidal2022hybrid}.
In the IRP-LR, there is a feasible alternative. The deliveries already on a
full vehicle often include quantities that are not needed to prevent future
stockouts, and this slack can be reassigned to the inserted customer without
violating any constraint. Exploiting this slack keeps the search within
capacity while giving access to the moves that matter most for the
logistic ratio.

Together, these observations identify the gap addressed in this paper:
there is no scalable heuristic for the IRP-LR that (i) evaluates routing and
quantity decisions jointly within each move, (ii) explores
inventory-infeasible solutions in a controlled manner that adapts to the
ratio objective, and (iii) does not depend on a mathematical programming
solver.

\begin{table*}[!htbp]
\centering
\caption{Studies most relevant to the IRP-LR and to the proposed solution
methodology.}
\label{tab:literature_irplr}
\scriptsize
\renewcommand{\arraystretch}{1.18}
\setlength{\tabcolsep}{3pt}
\begin{tabularx}{\textwidth}{
>{\raggedright\arraybackslash}p{2.3cm}
>{\raggedright\arraybackslash}p{2.6cm}
>{\raggedright\arraybackslash}p{2.8cm}
>{\raggedright\arraybackslash}X
>{\raggedright\arraybackslash}X}
\toprule
\textbf{Study} & \textbf{Problem / objective} & \textbf{Solution approach}
& \textbf{Key search characteristic} & \textbf{Relevance to this work} \\
\midrule
\citet{archetti2017logistic} & IRP; routing cost / quantity delivered
  & Exact fractional optimisation & Joint routing and replenishment
  & Introduces the IRP-LR \\
\citet{alvarez2018iterated} & IRP including logistic ratio
  & ILS and SA & Neighbourhood search over routes and replenishment
  & Benchmark heuristic for the IRP-LR \\
\citet{archetti2019exact} & IRP-LR & Iterative exact method
  & Sequence of IRPs with linear objectives
  & Optimal solutions for small instances \\
\citet{baldacci2020fractional} & VRP with fractional objective
  & Exact method & Primal and dual bounding procedures
  & General fractional routing methodology \\
\citet{absi2020branch} & Rich IRP; unit distribution cost
  & Heuristic branch-cut-and-price
  & Integrated routing, quantity, and scheduling
  & Large-scale industrial ratio IRP \\
\citet{su2020matheuristic} & ROADEF/EURO IRP & Local search with MIP/LP
  & Mathematical programming sets timing and quantities
  & Hybrid local-search design \\
\citet{he2020fixedsequence} & ROADEF/EURO IRP
  & Fixed-sequence matheuristic
  & Reoptimises quantities and timing for fixed visit sequences
  & Separation of route structure and quantity optimisation \\
\citet{alonsopecina2024iterated} & Operating cost / quantity delivered
  & ILS & Neighbourhood search for a multi-source IRP
  & Recent ILS with a ratio objective \\
\citet{man2025logistics} & Multi-product logistics-ratio IRP
  & Robust fractional optimisation & Uncertainty and lateral transshipment
  & Extension of logistics-ratio optimisation \\
\citet{castro2026hybrid} & Classical IRP & HGS
  & Route search with min-cost-flow quantity optimisation
  & State-of-the-art decomposition-based search \\
\citet{zhao2025large} & Classical and stockout-penalised IRP
  & HGS with large neighbourhood
  & Cross-period route/quantity moves; penalised stockouts
  & Joint moves and controlled infeasibility \\
\midrule
\textbf{This paper} & IRP-LR & Multi-start ILS (solver-free)
  & Joint route--quantity moves; vehicle slack; Lagrangian stockout penalty
  & Scalable heuristic for the IRP-LR \\
\bottomrule
\end{tabularx}
\end{table*}

\subsection{Contributions}
\label{sec:intro_contributions}

This paper proposes a multi-start iterated local search (MSILS) for the
IRP-LR. The algorithm does not require a mathematical programming solver.
Its main contributions are as follows.

\begin{enumerate}
  \item \textbf{Joint evaluation of route and quantity decisions.} Four of
  the five local search operators update delivery quantities as part of
  each move, using closed-form rules that respect depot supply, vehicle
  capacity, and customer storage capacity. Every candidate move is
  therefore evaluated by its actual effect on the logistic ratio, that is,
  on routing cost and delivered quantity together.

  \item \textbf{Vehicle slack.} We introduce a mechanism that reassigns
  quantity not needed to prevent future stockouts from the existing
  deliveries of a fully loaded vehicle to a newly inserted customer. The
  mechanism is motivated by the capacity-violation strategy of HGS
  \citep{vidal2022hybrid}, but it never violates vehicle capacity. It makes
  available the inter-route and inter-period moves that would otherwise be
  rejected, and it creates no new stockouts.

  \item \textbf{Lagrangian-guided exploration of infeasible solutions.}
  The zero-stockout constraint is dualised, and the local search solves a
  sequence of relaxed problems in which the Lagrange multiplier is updated
  by a subgradient-type rule based on the gap from the incumbent. The
  penalty therefore scales automatically to the magnitude of the logistic
  ratio of each instance, allowing the search to cross infeasible regions
  and then return to feasibility.

  \item \textbf{Problem-specific construction, rebalancing, and
  perturbation.} Initial solutions are built by a randomised constructive
  heuristic that prioritises customers by stockout urgency and distance
  from the depot. A rebalance procedure redistributes quantities over a
  fixed schedule whenever the incumbent improves, and an adaptive
  perturbation increases its strength as the number of non-improving
  iterations grows.

  \item \textbf{Computational evidence.} Experiments on
  \textcolor{red}{1128} benchmark instances, comprising \textcolor{red}{978}
  instances with 5 to 50 customers and \textcolor{red}{150} instances with
  50 to 200 customers, show that the MSILS improves on the ILS and SA of
  \citet{alvarez2018iterated} by on average \textcolor{red}{6.30\%} and
  \textcolor{red}{8.00\%}, respectively, on the large instances, with larger
  gains as the instance size grows. On small instances, it finds solutions
  within \textcolor{red}{X\%} of the optimal values reported by
  \citet{archetti2019exact}. A component analysis quantifies the
  contribution of each mechanism, and we also examine whether embedding the
  HGS of \citet{vidal2022hybrid} as a routing post-optimiser brings further
  improvement. The instances are publicly available
  \textcolor{red}{(and the code will be released upon acceptance)}.
\end{enumerate}

The remainder of the paper is organised as follows.
Section~\ref{sec:problem_description} defines the IRP-LR and the
Lagrangian-relaxed problem used by the algorithm. Section~\ref{sec: alg}
describes the MSILS. Section~\ref{sec: computational experiements} reports
the computational experiments, and Section~\textcolor{red}{\ref{sec:conclusion}}
concludes and outlines directions for future research.
 in place of the current
%        "Introduction", "Main Contributions" and
%        "Solution Approaches in the Literature" blocks.
% Requires: natbib (\citet/\citep), xcolor, tabularx, booktabs.
% Red text = numbers or claims to confirm before submission.
% =====================================================================

\section{Introduction}
\label{sec:introduction}

The inventory routing problem (IRP) arises in vendor-managed inventory
systems, in which a supplier monitors the stock levels of its customers and
takes responsibility for replenishing them. Over a multi-period planning
horizon, the supplier jointly decides when each customer is visited, how much
product is delivered at each visit, and how the vehicles are routed, subject
to customer storage capacities, vehicle capacities, and the requirement that
no customer runs out of stock. Integrating these decisions, rather than
fixing replenishment plans first and routing them afterwards, can produce
substantial savings \citep{archetti2016inventory}. However, the integration
couples a vehicle routing problem in every period with inventory dynamics
that link the periods together. As a result, the IRP remains one of the most
challenging problems in distribution logistics
\citep{moin2007overview,bertazzi2012inventory,coelho2014thirty,ekcsi2026inventory}.

Most of the IRP literature minimises the total cost, defined as the sum of
routing and inventory holding costs. In many applications, however, holding
costs are difficult to estimate or are negligible compared with transport
costs, and practitioners evaluate distribution plans by the transport cost
per unit of product delivered. This key performance indicator, known as the
\emph{logistic ratio}, has long been used in the distribution of industrial
gases \citep{bell1983improving,benoist2011randomized,singh2015incremental}.
It was also the objective of the ROADEF/EURO Challenge 2016, proposed by
Air Liquide \citep{andre2020roadef}. \citet{archetti2017logistic} formalised
the IRP with logistic ratio (IRP-LR) and showed that solutions that minimise
total cost can be poor in terms of the logistic ratio. One reason is the
finite-horizon effect: under a cost objective there is little incentive to
deliver product near the end of the horizon if it will only be consumed
afterwards. The logistic ratio, in contrast, rewards every unit delivered
and therefore produces replenishment plans that differ markedly from
cost-minimising ones.

The IRP-LR is harder to solve than the classical IRP. Its objective is the
ratio of two decision-dependent quantities: the routing cost in the numerator
and the total delivered quantity in the denominator. A move that shortens the
routes can therefore worsen the objective if it also reduces the quantity
delivered, and a longer route can be preferable if it allows a larger
delivery. Exact approaches build on parametric fractional programming
\citep{dinkelbach1967}, which replaces the ratio with a sequence of linear
objectives. \citet{archetti2017logistic} and \citet{archetti2019exact}
developed exact algorithms for the IRP-LR along these lines, and
\citet{baldacci2020fractional} proposed a general exact framework for vehicle
routing problems with fractional objectives. These methods provide optimal
solutions and bounds, but their computational requirements restrict them to
instances with at most \textcolor{red}{50} customers. For larger instances, to the best of our
knowledge, the only heuristics evaluated on the IRP-LR benchmark are the
iterated local search (ILS) and simulated annealing (SA) algorithms of
\citet{alvarez2018iterated}. The algorithms developed for the ROADEF/EURO
Challenge 2016 address a richer industrial variant with a ratio objective,
using heuristic branch-cut-and-price \citep{absi2020branch}, local search
combined with mathematical programming \citep{su2020matheuristic},
fixed-sequence reoptimisation of quantities and timing
\citep{he2020fixedsequence}, and hyper-heuristic sequence selection
\citep{kheiri2020heuristic}. More recently, ratio objectives have been
studied in a multi-source setting solved by ILS \citep{alonsopecina2024iterated}
and in a multi-product setting with lateral transshipment and uncertain
demand \citep{man2025logistics}. Compared with the classical IRP, for which
a wide range of exact and heuristic methods exists, the heuristic literature
on the IRP-LR remains sparse.

\subsection{Related literature and research gaps}
\label{sec:intro_literature}

For the classical IRP, exact methods based on branch-and-cut
\citep{archetti2007branch,coelho2013branch}, branch-price-and-cut
\citep{desaulniers2016branch}, and their matheuristic extensions
\citep{skaalnes2023branch} solve instances of moderate size. Large instances
are handled by heuristics, which are reviewed by \citet{sofianopoulou2021review}
and \citet{ekcsi2026inventory}. We organise the discussion around three
design choices that are particularly important under a logistic-ratio
objective. Table~\ref{tab:literature_irplr} summarises the studies most
closely related to this work.

\paragraph{Where delivery quantities are decided.}
Efficient local search is available for routing decisions, whereas delivery
quantities are naturally determined by mathematical programming. Most IRP
heuristics therefore decompose the problem: a heuristic explores visit
schedules and routes, and a mathematical program sets or repairs the
quantities \citep{iswari2026extending}. Examples include the tabu search
combined with mixed-integer programming of \citet{archetti2012hybrid}, the
matheuristic of \citet{archetti2017matheuristic}, the variable MIP
neighbourhood descent of \citet{coelho2020variable}, the ILS matheuristic of
\citet{lagana2024iterated}, and route-improving matheuristics for related
production routing problems \citep{vadseth2023multi}. In the hybrid genetic
search (HGS) of \citet{castro2026hybrid}, the routing component builds on the
HGS for the capacitated vehicle routing problem (CVRP)
\citep{vidal2022hybrid}, while the quantities are obtained by solving a
minimum-cost flow problem with the network simplex algorithm.
\citet{zhao2025large} extend HGS with a large neighbourhood that
simultaneously changes visit periods, insertion positions, and delivery
quantities across the horizon. Decomposition works well when holding costs
make quantity decisions secondary to routing. Under the logistic ratio,
however, the change in delivered quantity caused by a routing move is part
of the value of that move. A move that is evaluated with quantities held
fixed can appear unattractive even though it becomes improving once the
quantities are adjusted, and solving a mathematical program for every
candidate move is too expensive within a local search. This suggests
evaluating each move by its combined effect on routing cost and delivered
quantity. It also suggests using an exact optimisation step for quantities
only occasionally, as \citet{he2020fixedsequence} showed with fixed-sequence
reoptimisation.

\paragraph{How infeasibility is handled.}
Allowing a search to visit infeasible solutions and penalising the violation
in the objective is a well-established strategy in vehicle routing. It is a
central component of HGS \citep{vidal2022hybrid}, and in inventory routing
it has been used to relax vehicle capacity
\citep{archetti2012hybrid,castro2026hybrid} and customer stockouts
\citep{zhao2025large}. In the IRP, the stockout constraints make the
feasible region particularly fragmented. Removing a single visit usually
causes a stockout in some later period, so a search restricted to feasible
solutions cannot easily move between replenishment plans with different
visit patterns. Penalising stockouts addresses this, but under a ratio
objective the penalty weight is difficult to set. The logistic ratio is
measured in cost per unit and varies widely across instances, so a fixed
penalty is either too weak, and the search remains infeasible, or too
strong, and the search behaves as if infeasibility were forbidden. Lagrangian
relaxation \citep{fisher1981lagrangian} offers a principled alternative: the
stockout constraint is dualised, and the multiplier is updated with a
subgradient-type rule that scales the penalty to the current gap from the
incumbent. To the best of our knowledge, this has not previously been used
to guide a local search for the IRP-LR.

\paragraph{Saturated vehicles.}
Because the logistic ratio rewards large deliveries, vehicles in good
IRP-LR solutions are typically fully loaded. Any move that inserts a
customer into such a vehicle is then infeasible if the existing quantities
are kept, so large parts of the inter-route and inter-period neighbourhoods
are effectively unavailable. In the CVRP, HGS addresses the corresponding
issue by allowing capacity violations at a penalty \citep{vidal2022hybrid}.
In the IRP-LR, there is a feasible alternative. The deliveries already on a
full vehicle often include quantities that are not needed to prevent future
stockouts, and this slack can be reassigned to the inserted customer without
violating any constraint. Exploiting this slack keeps the search within
capacity while giving access to the moves that matter most for the
logistic ratio.

Together, these observations identify the gap addressed in this paper:
there is no scalable heuristic for the IRP-LR that (i) evaluates routing and
quantity decisions jointly within each move, (ii) explores
inventory-infeasible solutions in a controlled manner that adapts to the
ratio objective, and (iii) does not depend on a mathematical programming
solver.

\begin{table*}[!htbp]
\centering
\caption{Studies most relevant to the IRP-LR and to the proposed solution
methodology.}
\label{tab:literature_irplr}
\scriptsize
\renewcommand{\arraystretch}{1.18}
\setlength{\tabcolsep}{3pt}
\begin{tabularx}{\textwidth}{
>{\raggedright\arraybackslash}p{2.3cm}
>{\raggedright\arraybackslash}p{2.6cm}
>{\raggedright\arraybackslash}p{2.8cm}
>{\raggedright\arraybackslash}X
>{\raggedright\arraybackslash}X}
\toprule
\textbf{Study} & \textbf{Problem / objective} & \textbf{Solution approach}
& \textbf{Key search characteristic} & \textbf{Relevance to this work} \\
\midrule
\citet{archetti2017logistic} & IRP; routing cost / quantity delivered
  & Exact fractional optimisation & Joint routing and replenishment
  & Introduces the IRP-LR \\
\citet{alvarez2018iterated} & IRP including logistic ratio
  & ILS and SA & Neighbourhood search over routes and replenishment
  & Benchmark heuristic for the IRP-LR \\
\citet{archetti2019exact} & IRP-LR & Iterative exact method
  & Sequence of IRPs with linear objectives
  & Optimal solutions for small instances \\
\citet{baldacci2020fractional} & VRP with fractional objective
  & Exact method & Primal and dual bounding procedures
  & General fractional routing methodology \\
\citet{absi2020branch} & Rich IRP; unit distribution cost
  & Heuristic branch-cut-and-price
  & Integrated routing, quantity, and scheduling
  & Large-scale industrial ratio IRP \\
\citet{su2020matheuristic} & ROADEF/EURO IRP & Local search with MIP/LP
  & Mathematical programming sets timing and quantities
  & Hybrid local-search design \\
\citet{he2020fixedsequence} & ROADEF/EURO IRP
  & Fixed-sequence matheuristic
  & Reoptimises quantities and timing for fixed visit sequences
  & Separation of route structure and quantity optimisation \\
\citet{alonsopecina2024iterated} & Operating cost / quantity delivered
  & ILS & Neighbourhood search for a multi-source IRP
  & Recent ILS with a ratio objective \\
\citet{man2025logistics} & Multi-product logistics-ratio IRP
  & Robust fractional optimisation & Uncertainty and lateral transshipment
  & Extension of logistics-ratio optimisation \\
\citet{castro2026hybrid} & Classical IRP & HGS
  & Route search with min-cost-flow quantity optimisation
  & State-of-the-art decomposition-based search \\
\citet{zhao2025large} & Classical and stockout-penalised IRP
  & HGS with large neighbourhood
  & Cross-period route/quantity moves; penalised stockouts
  & Joint moves and controlled infeasibility \\
\midrule
\textbf{This paper} & IRP-LR & Multi-start ILS (solver-free)
  & Joint route--quantity moves; vehicle slack; Lagrangian stockout penalty
  & Scalable heuristic for the IRP-LR \\
\bottomrule
\end{tabularx}
\end{table*}

\subsection{Contributions}
\label{sec:intro_contributions}

This paper proposes a multi-start iterated local search (MSILS) for the
IRP-LR. The algorithm does not require a mathematical programming solver.
Its main contributions are as follows.

\begin{enumerate}
  \item \textbf{Joint evaluation of route and quantity decisions.} Four of
  the five local search operators update delivery quantities as part of
  each move, using closed-form rules that respect depot supply, vehicle
  capacity, and customer storage capacity. Every candidate move is
  therefore evaluated by its actual effect on the logistic ratio, that is,
  on routing cost and delivered quantity together.

  \item \textbf{Vehicle slack.} We introduce a mechanism that reassigns
  quantity not needed to prevent future stockouts from the existing
  deliveries of a fully loaded vehicle to a newly inserted customer. The
  mechanism is motivated by the capacity-violation strategy of HGS
  \citep{vidal2022hybrid}, but it never violates vehicle capacity. It makes
  available the inter-route and inter-period moves that would otherwise be
  rejected, and it creates no new stockouts.

  \item \textbf{Lagrangian-guided exploration of infeasible solutions.}
  The zero-stockout constraint is dualised, and the local search solves a
  sequence of relaxed problems in which the Lagrange multiplier is updated
  by a subgradient-type rule based on the gap from the incumbent. The
  penalty therefore scales automatically to the magnitude of the logistic
  ratio of each instance, allowing the search to cross infeasible regions
  and then return to feasibility.

  \item \textbf{Problem-specific construction, rebalancing, and
  perturbation.} Initial solutions are built by a randomised constructive
  heuristic that prioritises customers by stockout urgency and distance
  from the depot. A rebalance procedure redistributes quantities over a
  fixed schedule whenever the incumbent improves, and an adaptive
  perturbation increases its strength as the number of non-improving
  iterations grows.

  \item \textbf{Computational evidence.} Experiments on
  \textcolor{red}{1128} benchmark instances, comprising \textcolor{red}{978}
  instances with 5 to 50 customers and \textcolor{red}{150} instances with
  50 to 200 customers, show that the MSILS improves on the ILS and SA of
  \citet{alvarez2018iterated} by on average \textcolor{red}{6.30\%} and
  \textcolor{red}{8.00\%}, respectively, on the large instances, with larger
  gains as the instance size grows. On small instances, it finds solutions
  within \textcolor{red}{X\%} of the optimal values reported by
  \citet{archetti2019exact}. A component analysis quantifies the
  contribution of each mechanism, and we also examine whether embedding the
  HGS of \citet{vidal2022hybrid} as a routing post-optimiser brings further
  improvement. The instances are publicly available
  \textcolor{red}{(and the code will be released upon acceptance)}.
\end{enumerate}

The remainder of the paper is organised as follows.
Section~\ref{sec:problem_description} defines the IRP-LR and the
Lagrangian-relaxed problem used by the algorithm. Section~\ref{sec: alg}
describes the MSILS. Section~\ref{sec: computational experiements} reports
the computational experiments, and Section~\textcolor{red}{\ref{sec:conclusion}}
concludes and outlines directions for future research.
 in place of the current
%        "Introduction", "Main Contributions" and
%        "Solution Approaches in the Literature" blocks.
% Requires: natbib (\citet/\citep), xcolor, tabularx, booktabs.
% Red text = numbers or claims to confirm before submission.
% =====================================================================

\section{Introduction}
\label{sec:introduction}

The inventory routing problem (IRP) arises in vendor-managed inventory
systems, in which a supplier monitors the stock levels of its customers and
takes responsibility for replenishing them. Over a multi-period planning
horizon, the supplier jointly decides when each customer is visited, how much
product is delivered at each visit, and how the vehicles are routed, subject
to customer storage capacities, vehicle capacities, and the requirement that
no customer runs out of stock. Integrating these decisions, rather than
fixing replenishment plans first and routing them afterwards, can produce
substantial savings \citep{archetti2016inventory}. However, the integration
couples a vehicle routing problem in every period with inventory dynamics
that link the periods together. As a result, the IRP remains one of the most
challenging problems in distribution logistics
\citep{moin2007overview,bertazzi2012inventory,coelho2014thirty,ekcsi2026inventory}.

Most of the IRP literature minimises the total cost, defined as the sum of
routing and inventory holding costs. In many applications, however, holding
costs are difficult to estimate or are negligible compared with transport
costs, and practitioners evaluate distribution plans by the transport cost
per unit of product delivered. This key performance indicator, known as the
\emph{logistic ratio}, has long been used in the distribution of industrial
gases \citep{bell1983improving,benoist2011randomized,singh2015incremental}.
It was also the objective of the ROADEF/EURO Challenge 2016, proposed by
Air Liquide \citep{andre2020roadef}. \citet{archetti2017logistic} formalised
the IRP with logistic ratio (IRP-LR) and showed that solutions that minimise
total cost can be poor in terms of the logistic ratio. One reason is the
finite-horizon effect: under a cost objective there is little incentive to
deliver product near the end of the horizon if it will only be consumed
afterwards. The logistic ratio, in contrast, rewards every unit delivered
and therefore produces replenishment plans that differ markedly from
cost-minimising ones.

The IRP-LR is harder to solve than the classical IRP. Its objective is the
ratio of two decision-dependent quantities: the routing cost in the numerator
and the total delivered quantity in the denominator. A move that shortens the
routes can therefore worsen the objective if it also reduces the quantity
delivered, and a longer route can be preferable if it allows a larger
delivery. Exact approaches build on parametric fractional programming
\citep{dinkelbach1967}, which replaces the ratio with a sequence of linear
objectives. \citet{archetti2017logistic} and \citet{archetti2019exact}
developed exact algorithms for the IRP-LR along these lines, and
\citet{baldacci2020fractional} proposed a general exact framework for vehicle
routing problems with fractional objectives. These methods provide optimal
solutions and bounds, but their computational requirements restrict them to
instances with at most \textcolor{red}{50} customers. For larger instances, to the best of our
knowledge, the only heuristics evaluated on the IRP-LR benchmark are the
iterated local search (ILS) and simulated annealing (SA) algorithms of
\citet{alvarez2018iterated}. The algorithms developed for the ROADEF/EURO
Challenge 2016 address a richer industrial variant with a ratio objective,
using heuristic branch-cut-and-price \citep{absi2020branch}, local search
combined with mathematical programming \citep{su2020matheuristic},
fixed-sequence reoptimisation of quantities and timing
\citep{he2020fixedsequence}, and hyper-heuristic sequence selection
\citep{kheiri2020heuristic}. More recently, ratio objectives have been
studied in a multi-source setting solved by ILS \citep{alonsopecina2024iterated}
and in a multi-product setting with lateral transshipment and uncertain
demand \citep{man2025logistics}. Compared with the classical IRP, for which
a wide range of exact and heuristic methods exists, the heuristic literature
on the IRP-LR remains sparse.

\subsection{Related literature and research gaps}
\label{sec:intro_literature}

For the classical IRP, exact methods based on branch-and-cut
\citep{archetti2007branch,coelho2013branch}, branch-price-and-cut
\citep{desaulniers2016branch}, and their matheuristic extensions
\citep{skaalnes2023branch} solve instances of moderate size. Large instances
are handled by heuristics, which are reviewed by \citet{sofianopoulou2021review}
and \citet{ekcsi2026inventory}. We organise the discussion around three
design choices that are particularly important under a logistic-ratio
objective. Table~\ref{tab:literature_irplr} summarises the studies most
closely related to this work.

\paragraph{Where delivery quantities are decided.}
Efficient local search is available for routing decisions, whereas delivery
quantities are naturally determined by mathematical programming. Most IRP
heuristics therefore decompose the problem: a heuristic explores visit
schedules and routes, and a mathematical program sets or repairs the
quantities \citep{iswari2026extending}. Examples include the tabu search
combined with mixed-integer programming of \citet{archetti2012hybrid}, the
matheuristic of \citet{archetti2017matheuristic}, the variable MIP
neighbourhood descent of \citet{coelho2020variable}, the ILS matheuristic of
\citet{lagana2024iterated}, and route-improving matheuristics for related
production routing problems \citep{vadseth2023multi}. In the hybrid genetic
search (HGS) of \citet{castro2026hybrid}, the routing component builds on the
HGS for the capacitated vehicle routing problem (CVRP)
\citep{vidal2022hybrid}, while the quantities are obtained by solving a
minimum-cost flow problem with the network simplex algorithm.
\citet{zhao2025large} extend HGS with a large neighbourhood that
simultaneously changes visit periods, insertion positions, and delivery
quantities across the horizon. Decomposition works well when holding costs
make quantity decisions secondary to routing. Under the logistic ratio,
however, the change in delivered quantity caused by a routing move is part
of the value of that move. A move that is evaluated with quantities held
fixed can appear unattractive even though it becomes improving once the
quantities are adjusted, and solving a mathematical program for every
candidate move is too expensive within a local search. This suggests
evaluating each move by its combined effect on routing cost and delivered
quantity. It also suggests using an exact optimisation step for quantities
only occasionally, as \citet{he2020fixedsequence} showed with fixed-sequence
reoptimisation.

\paragraph{How infeasibility is handled.}
Allowing a search to visit infeasible solutions and penalising the violation
in the objective is a well-established strategy in vehicle routing. It is a
central component of HGS \citep{vidal2022hybrid}, and in inventory routing
it has been used to relax vehicle capacity
\citep{archetti2012hybrid,castro2026hybrid} and customer stockouts
\citep{zhao2025large}. In the IRP, the stockout constraints make the
feasible region particularly fragmented. Removing a single visit usually
causes a stockout in some later period, so a search restricted to feasible
solutions cannot easily move between replenishment plans with different
visit patterns. Penalising stockouts addresses this, but under a ratio
objective the penalty weight is difficult to set. The logistic ratio is
measured in cost per unit and varies widely across instances, so a fixed
penalty is either too weak, and the search remains infeasible, or too
strong, and the search behaves as if infeasibility were forbidden. Lagrangian
relaxation \citep{fisher1981lagrangian} offers a principled alternative: the
stockout constraint is dualised, and the multiplier is updated with a
subgradient-type rule that scales the penalty to the current gap from the
incumbent. To the best of our knowledge, this has not previously been used
to guide a local search for the IRP-LR.

\paragraph{Saturated vehicles.}
Because the logistic ratio rewards large deliveries, vehicles in good
IRP-LR solutions are typically fully loaded. Any move that inserts a
customer into such a vehicle is then infeasible if the existing quantities
are kept, so large parts of the inter-route and inter-period neighbourhoods
are effectively unavailable. In the CVRP, HGS addresses the corresponding
issue by allowing capacity violations at a penalty \citep{vidal2022hybrid}.
In the IRP-LR, there is a feasible alternative. The deliveries already on a
full vehicle often include quantities that are not needed to prevent future
stockouts, and this slack can be reassigned to the inserted customer without
violating any constraint. Exploiting this slack keeps the search within
capacity while giving access to the moves that matter most for the
logistic ratio.

Together, these observations identify the gap addressed in this paper:
there is no scalable heuristic for the IRP-LR that (i) evaluates routing and
quantity decisions jointly within each move, (ii) explores
inventory-infeasible solutions in a controlled manner that adapts to the
ratio objective, and (iii) does not depend on a mathematical programming
solver.

\begin{table*}[!htbp]
\centering
\caption{Studies most relevant to the IRP-LR and to the proposed solution
methodology.}
\label{tab:literature_irplr}
\scriptsize
\renewcommand{\arraystretch}{1.18}
\setlength{\tabcolsep}{3pt}
\begin{tabularx}{\textwidth}{
>{\raggedright\arraybackslash}p{2.3cm}
>{\raggedright\arraybackslash}p{2.6cm}
>{\raggedright\arraybackslash}p{2.8cm}
>{\raggedright\arraybackslash}X
>{\raggedright\arraybackslash}X}
\toprule
\textbf{Study} & \textbf{Problem / objective} & \textbf{Solution approach}
& \textbf{Key search characteristic} & \textbf{Relevance to this work} \\
\midrule
\citet{archetti2017logistic} & IRP; routing cost / quantity delivered
  & Exact fractional optimisation & Joint routing and replenishment
  & Introduces the IRP-LR \\
\citet{alvarez2018iterated} & IRP including logistic ratio
  & ILS and SA & Neighbourhood search over routes and replenishment
  & Benchmark heuristic for the IRP-LR \\
\citet{archetti2019exact} & IRP-LR & Iterative exact method
  & Sequence of IRPs with linear objectives
  & Optimal solutions for small instances \\
\citet{baldacci2020fractional} & VRP with fractional objective
  & Exact method & Primal and dual bounding procedures
  & General fractional routing methodology \\
\citet{absi2020branch} & Rich IRP; unit distribution cost
  & Heuristic branch-cut-and-price
  & Integrated routing, quantity, and scheduling
  & Large-scale industrial ratio IRP \\
\citet{su2020matheuristic} & ROADEF/EURO IRP & Local search with MIP/LP
  & Mathematical programming sets timing and quantities
  & Hybrid local-search design \\
\citet{he2020fixedsequence} & ROADEF/EURO IRP
  & Fixed-sequence matheuristic
  & Reoptimises quantities and timing for fixed visit sequences
  & Separation of route structure and quantity optimisation \\
\citet{alonsopecina2024iterated} & Operating cost / quantity delivered
  & ILS & Neighbourhood search for a multi-source IRP
  & Recent ILS with a ratio objective \\
\citet{man2025logistics} & Multi-product logistics-ratio IRP
  & Robust fractional optimisation & Uncertainty and lateral transshipment
  & Extension of logistics-ratio optimisation \\
\citet{castro2026hybrid} & Classical IRP & HGS
  & Route search with min-cost-flow quantity optimisation
  & State-of-the-art decomposition-based search \\
\citet{zhao2025large} & Classical and stockout-penalised IRP
  & HGS with large neighbourhood
  & Cross-period route/quantity moves; penalised stockouts
  & Joint moves and controlled infeasibility \\
\midrule
\textbf{This paper} & IRP-LR & Multi-start ILS (solver-free)
  & Joint route--quantity moves; vehicle slack; Lagrangian stockout penalty
  & Scalable heuristic for the IRP-LR \\
\bottomrule
\end{tabularx}
\end{table*}

\subsection{Contributions}
\label{sec:intro_contributions}

This paper proposes a multi-start iterated local search (MSILS) for the
IRP-LR. The algorithm does not require a mathematical programming solver.
Its main contributions are as follows.

\begin{enumerate}
  \item \textbf{Joint evaluation of route and quantity decisions.} Four of
  the five local search operators update delivery quantities as part of
  each move, using closed-form rules that respect depot supply, vehicle
  capacity, and customer storage capacity. Every candidate move is
  therefore evaluated by its actual effect on the logistic ratio, that is,
  on routing cost and delivered quantity together.

  \item \textbf{Vehicle slack.} We introduce a mechanism that reassigns
  quantity not needed to prevent future stockouts from the existing
  deliveries of a fully loaded vehicle to a newly inserted customer. The
  mechanism is motivated by the capacity-violation strategy of HGS
  \citep{vidal2022hybrid}, but it never violates vehicle capacity. It makes
  available the inter-route and inter-period moves that would otherwise be
  rejected, and it creates no new stockouts.

  \item \textbf{Lagrangian-guided exploration of infeasible solutions.}
  The zero-stockout constraint is dualised, and the local search solves a
  sequence of relaxed problems in which the Lagrange multiplier is updated
  by a subgradient-type rule based on the gap from the incumbent. The
  penalty therefore scales automatically to the magnitude of the logistic
  ratio of each instance, allowing the search to cross infeasible regions
  and then return to feasibility.

  \item \textbf{Problem-specific construction, rebalancing, and
  perturbation.} Initial solutions are built by a randomised constructive
  heuristic that prioritises customers by stockout urgency and distance
  from the depot. A rebalance procedure redistributes quantities over a
  fixed schedule whenever the incumbent improves, and an adaptive
  perturbation increases its strength as the number of non-improving
  iterations grows.

  \item \textbf{Computational evidence.} Experiments on
  \textcolor{red}{1128} benchmark instances, comprising \textcolor{red}{978}
  instances with 5 to 50 customers and \textcolor{red}{150} instances with
  50 to 200 customers, show that the MSILS improves on the ILS and SA of
  \citet{alvarez2018iterated} by on average \textcolor{red}{6.30\%} and
  \textcolor{red}{8.00\%}, respectively, on the large instances, with larger
  gains as the instance size grows. On small instances, it finds solutions
  within \textcolor{red}{X\%} of the optimal values reported by
  \citet{archetti2019exact}. A component analysis quantifies the
  contribution of each mechanism, and we also examine whether embedding the
  HGS of \citet{vidal2022hybrid} as a routing post-optimiser brings further
  improvement. The instances are publicly available
  \textcolor{red}{(and the code will be released upon acceptance)}.
\end{enumerate}

The remainder of the paper is organised as follows.
Section~\ref{sec:problem_description} defines the IRP-LR and the
Lagrangian-relaxed problem used by the algorithm. Section~\ref{sec: alg}
describes the MSILS. Section~\ref{sec: computational experiements} reports
the computational experiments, and Section~\textcolor{red}{\ref{sec:conclusion}}
concludes and outlines directions for future research.
